% MT032 / explanation — canonical editable content.
Bayes分类比较的是看到\(x\)后，各类的后验概率。错分代价相同时，选概率最大的类最有把握：若两类后验为0.7和0.3，选第一类的错误概率为0.3，已经比选第二类的0.7小。



\begin{itemize}

\item \textbf{A错误。} NB的边界取决于类条件分布。Gaussian NB的类间方差不同，比较分数时可以留下平方项，形成非线性边界。

\item \textbf{B依本卷语境不选。} 生成式模型直接规定\(p(x\mid y)\)和\(P(y)\)，再由Bayes公式推得\(P(y\mid x)\)。若“显式定义”只是指能写出后验计算式，字面上又可成立；原卷按“直接建模的对象”区分生成式与判别式，配套解答只选D。

\item \textbf{C错误。} 有限数据估出的分布也可能过度贴合样本。比如只见到一次某个事件就把概率估成1，未必适合新数据。

\item \textbf{D按原解选。} 使用正确后验、错分代价相同的Bayes决策，最小化每个输入位置的错误概率，因此也最小化平均错误率。拟合出的NB若概率估计错了，不能自动享有这个保证。

\item \textbf{E错误。} 后验分数包含类别先验。某类本来极罕见，即使新输入略像它，也可能不足以抵消先验差异；数值例子见\StudyRef{Answers}{mt040}{MT040}。

\end{itemize}


作答为D，并能解释B的语义。这里的歧义来自原题措辞，同学笔记中的BD与配套答案D应保留区分。


\textbf{出处：} 2021B*期中 Q6，原卷第3页（5分）


\textbf{答案依据：} 2021B配套解答 / 2021B配套答案副本 Q6答案D；同学小抄写BD，保留措辞解释。
